\section{Spectral transfer and the bound at every scale}
\label{q:sec-spectral}

The moment bound for $H$ has only square-root dependence on each
prime supply's reciprocal mass. We now use it to control the retained
centered sums $C_j^\circ$. First we transfer control of $H$ to the
projected edge operator $E(\sum_d A_d)E$, then test against the
Liouville function with vertex weight $\sqrt w$. The factor $L$ in
the edge matrices yields a factor $1/L$ in the estimate for each bin,
compensating for the factor $L$ in the number $O(L/\eta)$ of bins.
After division by the retained mass $S_0\ge SV_*/2$, with $V_*\ge W^J$,
the spectral contribution will have the form
\[
 \frac K\eta\left(\frac C{\sqrt W}\right)^J
\]
for an absolute constant $C$. We prove the transfer and the required
finite-interval estimates before choosing $W$.

The use of a nonbacktracking operator to control an adjacency operator
is exemplified by the weighted Ihara--Bass formula in
\cite[Section~4.2]{Pilatte26}.  We prove the transfer needed here
directly for self-adjoint edge operators; distinct edge operators need
not commute.

\subsection{A transfer lemma for self-adjoint edge matrices}

\begin{lemma}\label{q:noncommuting}
Let $N\ge1$ and let $B_1,\ldots,B_N$ be self-adjoint operators on a finite-dimensional
complex Hilbert space $\mathcal H$, and let $P$ be an orthogonal
projection on $\mathcal H$. Define an operator $\mathscr B$ on
$\mathcal H^N$ by
\[
 \mathscr B_{i,j}=\1_{i\ne j}B_j.
\]
Suppose $a,b\ge0$, $\|B_i\|\le a$ for every $i$, and
\[
 \sum_{i=1}^N\|B_i v\|^2\le b^2\|v\|^2
 \qquad(v=Pv).
\]
Writing $\rho(\mathscr B)$ for its spectral radius, one has
\begin{equation}\label{q:noncommuting-bound}
 \left\|P\left(\sum_{i=1}^N B_i\right)P\right\|
 \le 3\max\{a,b,\rho(\mathscr B)\}.
\end{equation}
\end{lemma}

\begin{proof}
Put $m=\max\{a,b,\rho(\mathscr B)\}$. If $m=0$, all $B_i$ vanish.
Otherwise set $t=(2m)^{-1}$. For real $|u|\le t$, the operators
$I-u\mathscr B$ and $I+uB_i$ are invertible. Define
\[
 F(u)=I-\sum_{i=1}^N uB_i(I+uB_i)^{-1}.
\]
Each summand is self-adjoint, since $B_i$ commutes with its own
resolvent. If $F(u)v=0$, set $z_i=(I+uB_i)^{-1}v$. Then
\[
 (u\mathscr Bz)_i
 =\sum_{j\ne i}uB_jz_j
 =v-uB_i z_i=z_i.
\]
A nonzero $v$ would give a nonzero $z$, contradicting the invertibility
of $I-u\mathscr B$. Hence $F(u)$ is invertible throughout $[-t,t]$.
It is positive definite there by continuity, because $F(0)=I$.

The resolvent identity gives
\[
 F(u)=I-u\sum_i B_i+Q(u),\qquad
 Q(u)=u^2\sum_i B_i^2(I+uB_i)^{-1}.
\]
Since $|u|\|B_i\|\le1/2$, the inverses in the last expression are
positive and bounded above by $2I$. For $v=Pv$,
\[
 0\le\langle v,Q(u)v\rangle
 \le2u^2\sum_i\|B_i v\|^2
 \le2u^2b^2\|v\|^2.
\]
Positivity of $F(t)$ and $F(-t)$ therefore yields
\[
 \left|\left\langle v,\sum_i B_i v\right\rangle\right|
 \le \frac{1+2t^2b^2}{t}\|v\|^2
 \le3m\|v\|^2.
\]
Taking the supremum over unit vectors in the range of $P$ proves
\eqref{q:noncommuting-bound}. No step commutes two distinct $B_i$.
\end{proof}

\subsection{Weighted row bounds and the degree projection}

Fix a bin $j$ and one block of $M$ sites, with its matrices $A_d$, $H$
and projection $E$ from the trace construction. Define
\[
 a_d(n)=\prod_{p\mid d}\left|\1_{p\mid n}-\frac1p\right|.
\]
If $A_d(n,n')\ne0$, the difference $n'-n$ is a multiple of $d$.
Thus $a_d(n)=a_d(n')$ along every edge of $A_d$.

\begin{lemma}\label{q:weighted-rows}
On every block, in either the product model or the integer model,
\begin{align}
 \|A_d\|&\le2K,\label{q:single-operator}\\
 \sum_{d\in\mathcal D}\|A_d v\|^2
 &\le4K^2(8W)^J\|v\|^2\qquad(v=Ev).
 \label{q:square-operators}
\end{align}
\end{lemma}

\begin{proof}
For each of the two edge orientations, the weighted absolute row sum
at $n$ is at most
\[
 a_d(n)\sum_q\frac{Lu(q)}{w(n)}\1_{q\mid n}
 \1_{(d,q)\in\mathcal S_j}\1_{\rho_j(n,d)\le K/L}
 \le K a_d(n).
\]
Indeed the sum without the last indicator is at most
$L\rho_j(n,d)$. For an incoming edge, $q$ divides one endpoint if
and only if it divides the other. All the other edge restrictions can
be discarded in this nonnegative bound. Consequently
\begin{equation}\label{q:weighted-row}
 \sum_{n'}|A_d(n,n')|
       \frac{\sqrt{w(n')}}{\sqrt{w(n)}}\le2K a_d(n).
\end{equation}

The weighted Schur inequality for a real symmetric matrix $B$ follows
by applying
\[
 2|v_n v_{n'}|\le
 |v_n|^2\frac{\sqrt{w(n')}}{\sqrt{w(n)}}+
 |v_{n'}|^2\frac{\sqrt{w(n)}}{\sqrt{w(n')}}
\]
to its quadratic form. Apply it to $A_d$ on each level set of $a_d$;
these level sets are invariant under $A_d$. The resulting bound is
\begin{equation}\label{q:localized-square}
 \|A_d v\|^2\le4K^2\sum_n a_d(n)^2|v_n|^2.
\end{equation}
In particular, $a_d(n)\le1$ proves \eqref{q:single-operator}.

At a site retained by $E$, the total P-degree is at most $6WJ$.
Since $V_i\le2W$, the arithmetic-geometric mean inequality gives
\begin{align*}
 \sum_d a_d(n)^2
 &\le\sum_d a_d(n)
 \le\prod_{i=1}^J\bigl(\omega_{P_i}(n)+V_i\bigr)\\
 &\le\left(\frac{\omega_P(n)+\sum_iV_i}{J}\right)^J
 \le(8W)^J.
\end{align*}
Sum \eqref{q:localized-square} over $d$. This proves
\eqref{q:square-operators} without requiring $E$ to commute with $A_d$.
\end{proof}

\begin{proposition}\label{q:spectral}
There is an absolute constant $C_3$ with the following property.
For each bin $j$, average block origins over any integer interval to
which \Cref{q:trace} applies. Outside a fraction at most $e^{-2k}$ of
these origins,
\begin{equation}\label{q:compressed-norm}
 \left\|E\left(\sum_d A_d\right)E\right\|
 \le K(C_3\sqrt W)^J.
\end{equation}
\end{proposition}

\begin{proof}
For every finite matrix, $\rho(H)^{2k}\le\|H^k\|_{\mathrm{HS}}^2$.
Thus \Cref{q:trace} and Markov's inequality show that
\[
 \rho(H)\le eK(C_2\sqrt W)^J
\]
except on a fraction at most $e^{-2k}$ of the origins. On each remaining
block apply \Cref{q:noncommuting} with
\[
 a=2K,\qquad b=2K(8W)^{J/2},\qquad P=E.
\]
For $J\ge1$ and $W\ge10$, the choice
\begin{equation}\label{q:C3-choice}
 C_3=3\max\{2,2\sqrt8,eC_2\}
\end{equation}
gives \eqref{q:compressed-norm}. In particular, $C_3$ is independent of
$W,h,l$.
\end{proof}

\subsection{Testing and averaging overlapping blocks}

We now turn \eqref{q:compressed-norm} into an estimate for $C_j^\circ$,
the retained correlation in \Cref{q:cuts}. Its endpoint restrictions
are exactly those imposed by the $A_d$ and the projection $E$.

\begin{proposition}\label{q:testing}
For every bin $j$ and all sufficiently large $X$,
\begin{equation}\label{q:bin-estimate}
 |C_j^\circ|\ll \frac{S K(C_3\sqrt W)^J}{L}+e^{-L}.
\end{equation}
The threshold may depend on the fixed $h,l,W$, while $C_3$ is the
absolute constant in \eqref{q:C3-choice}.
\end{proposition}

\begin{proof}
Write $T=T_j$ and $N=\lfloor T\rfloor$, and average over blocks
$I_t=\{t+1,\ldots,t+M\}$, $0\le t<N$. For large $X$,
$N\ge\exp(L^A/2)$, so the interval of origins is admissible in
\Cref{q:finite-law,q:trace}. On each block put
\[
 f_t(n)=\sqrt{w(n)}\lambda(n)
             \1_{\omega_Q(n)\le400\log L},\qquad g_t=Ef_t.
\]
We first check the averaged norm needed for testing:
\begin{equation}\label{q:test-vector-norm}
 \frac1N\sum_{t=0}^{N-1}\|g_t\|^2
 \le\frac1N\sum_{t=0}^{N-1}\|f_t\|^2\le2MS.
\end{equation}
For each of the $M$ block positions, \Cref{q:truncated-weight}
applied to the corresponding translated interval of $N$ origins
bounds the average truncated weight by $2S$. Summing proves
\eqref{q:test-vector-norm}. This uses the truncated comparison
already established, without an untruncated integer moment of $w$.

On blocks satisfying \eqref{q:compressed-norm}, the averaged absolute
quadratic form is at most $2MSK(C_3\sqrt W)^J$. On every block,
including exceptional ones, \eqref{q:weighted-row} gives the
deterministic bound
\begin{equation}\label{q:deterministic-test}
 \left|\left\langle g_t,\sum_d A_d g_t\right\rangle\right|
 \le 2KM\,5^{400\log L}(8W)^J.
\end{equation}
Here we bound the form by its entrywise absolute value, use
$|g_t(n)|\le\sqrt{w(n)}$, and then sum the weighted row bounds at
sites satisfying both degree cutoffs. The right side divided by $M$
is a fixed power of $L$ for fixed $W$. Multiplication by the exceptional
fraction $e^{-2k}$ makes its contribution, after division by $2LM$,
at most $e^{-L}$ for large $L$. Therefore
\begin{equation}\label{q:averaged-form}
 \left|\frac1{2LMN}\sum_{t=0}^{N-1}
       \left\langle g_t,\sum_d A_d g_t\right\rangle\right|
 \ll \frac{S K(C_3\sqrt W)^J}{L}+e^{-L}.
\end{equation}

It remains to compare this form with the original prefix sum. Denote
by $c_j(n;d,q)$ the signed summand belonging to $(d,q)$ in
$T C_j^\circ$, including all its endpoint cutoffs. Use the same
formula to define it for every positive $n$, and put $r=hqd$. These cutoffs
are predicates of the ambient sites, independent of the choice of
block; restricting to a block only removes edges leaving it.
Testing an increasing matrix entry cancels its two square-root
weights and gives $L c_j(n;d,q)$; symmetry gives the same term in the
opposite orientation. Thus the form on the left of
\eqref{q:averaged-form}, before taking absolute values, is exactly
\begin{equation}\label{q:edge-multiplicity}
 \frac1{MN}\sum_{n,d,q}m_N(n,r)c_j(n;d,q),\qquad
 m_N(n,r)=
 \bigl[\min(N-1,n-1)-\max(0,n+r-M)+1\bigr]_+,
\end{equation}
where $[x]_+=\max(x,0)$. This is the number of blocks containing both
endpoints. All contributing $n$ are positive.

Let
\[
 R_L=h e^{100L+1},\qquad
 D_L=\frac KL\,5^{400\log L}(8W)^J.
\]
Every displacement is at most $R_L$, and the one-orientation row
bound proves $\sum_{d,q}|c_j(n;d,q)|\le D_L$ for every $n$.
For large $L$ we have $N\gg M>R_L$. If $M\le n\le N$, then
$m_N(n,r)=M-r$. The lower boundary and the extra sites $N<n<N+M$
contain at most $2M$ sites, and always $0\le m_N(n,r)\le M$.
Consequently the difference between \eqref{q:edge-multiplicity} and
$C_j^\circ=T^{-1}\sum_{n\le N,d,q}c_j(n;d,q)$ is at most
\begin{equation}\label{q:prefix-error}
 O\left(D_L\left(\frac{R_L}{M}+\frac MN+\frac1T\right)\right).
\end{equation}
The last term accounts for replacing $N^{-1}$ by $T^{-1}$.
Since $M=\lceil e^{103L}\rceil$, $T\ge X=e^{L^A}$ and $D_L$ is
polynomial in $L$, \eqref{q:prefix-error} is $O(e^{-L})$ for
sufficiently large $L$. This also bounds edges whose terminal endpoint
exceeds $T$ and the floor errors. Combining with
\eqref{q:averaged-form} proves \eqref{q:bin-estimate}.
\end{proof}

\subsection{Final choice of constants}

\begin{proof}[Proof of \Cref{q:progression}]
There are $O(L/\eta)$ bins. Sum \eqref{q:bin-estimate}, use
$S_0\ge SV_*/2$ and $V_*\ge W^J$, and then apply the centering and
deletion estimates of \Cref{q:center,q:cuts}. We obtain
\begin{align}
 |F(X)|\ll_{h,l,W}{}&
 \eta+\eta^{-1}2^J L^{-1/20}
 +2^J\bigl(K^{-1}+L^{-100}+e^{-2WJ}\bigr)\notag\\
 &+\frac K\eta\left(\frac{C_3}{\sqrt W}\right)^J
 +e^{-L^{0.9}}+X^{-1}.
 \label{q:finish}
\end{align}
The deletion error in \Cref{q:cuts} is already summed over the bins.
The new error $O((L/\eta)e^{-L})$ from \Cref{q:testing} is absorbed
by $e^{-L^{0.9}}$, because $L/\eta$ is polynomial in $L$ and $SV_*\ge1$.

Choose the absolute constants in the following order. First fix $A$
as required by \Cref{q:finite-law}; in particular $A\ge1000$ and
$A>C_1+20$, where $C_1$ is the absolute exponent in that comparison.
The trace proof gives an absolute $C_2$, and
\eqref{q:C3-choice} then gives an absolute $C_3$. Now choose $W\ge10$
so large that
\begin{equation}\label{q:W-choice}
 e^5 C_3/\sqrt W\le e^{-1}.
\end{equation}
Only after fixing these constants do we increase the threshold for
$X$, allowing it to depend on $h,l,W$.

Recall $\eta=e^{-J}$, $K=e^{4J}$ and
$J\le\delta\log L/(6W)$ with $\delta=1/200$.
The spectral term in \eqref{q:finish} is at most $e^{-J}$ by
\eqref{q:W-choice}. The other terms satisfy
\begin{align*}
 2^JK^{-1}&=e^{-(4-\log2)J}\le e^{-J},\\
 2^Je^{-2WJ}&=e^{-(2W-\log2)J}\le e^{-J},\\
 e^J2^J L^{-1/20}&\le e^{-J},\qquad
 2^J L^{-100}\le e^{-J}.
\end{align*}
For the last line, it is enough to note that
\[
 \frac{\delta(2+\log2)}{6W}<\frac1{20},\qquad
 \frac{\delta(1+\log2)}{6W}<100.
\]
Also $e^{-L^{0.9}}+X^{-1}\ll e^{-J}$ for large $X$. Hence
$|F(X)|\ll_{h,l}e^{-J}$. Finally,
\[
 e^{-J}\le e L^{-\delta/(6W)}
 =e(\log X)^{-c},\qquad c=\frac{\delta}{6WA}>0.
\]
All three constants $A,W,c$ are absolute. The argument applies to
every sufficiently large real $X$, with no excluded scales. On the
remaining bounded range $3\le X\le X_0(h,l)$, the estimate follows
from $|F(X)|\le1$ after enlarging the implied constant. This proves
\Cref{q:progression}.
\end{proof}
